\section{Problem Formulation}

\subsection{Agent and target dynamics}

Consider $N$ fully-actuated agents and one moving target, all evolving in
$\RR^n$ ($n=2$ or $3$). Every agent is modelled as a double integrator
\begin{equation}
    \ddot p_i = u_i + d_i, \qquad i\in\NN=\{1,\dots,N\},
    \label{eq:agent}
\end{equation}
where $p_i\in\RR^n$ is the position, $u_i\in\RR^n$ the control input, and
$d_i$ a bounded disturbance. The target obeys the same dynamics
\begin{equation}
    \ddot p_0 = u_0 + d_0,
    \label{eq:target}
\end{equation}
with $u_0$ (and $d_0$) possibly unknown to the agents. Introduce the
relative-to-target states
\[
p_{i0}=p_i-p_0,\qquad v_{i0}=\dot p_i-\dot p_0.
\]

A natural measure of how well the swarm encloses the target is the
\emph{centroid error}
\begin{equation}
    \tilde p = \bar p - p_0
    = \frac{1}{N}\sum_{i\in\NN} p_i - p_0,
    \qquad \dot{\tilde p} = \frac{1}{N}\sum_{i\in\NN} v_i - v_0.
    \label{eq:centroid}
\end{equation}
Since the centroid $\bar p$ always belongs to the agents' convex hull
$\co(p)$, bounding $\tilde p$ bounds the target-to-hull distance (see
\eqref{eq:dist} below). Differentiating \eqref{eq:centroid} and substituting
\eqref{eq:agent}--\eqref{eq:target} gives the open-loop centroid dynamics
\begin{equation}
    \ddot{\tilde p}
    = \frac{1}{N}\sum_{i\in\NN}(u_i+d_i) - (u_0+d_0).
    \label{eq:centroid-dyn}
\end{equation}

\subsection{Fencing properties}

\begin{definition}[Fencing properties]
    (Adapted from \cite{chenCooperativeTargetfencingProtocol2019})
    The self-organizing fencing problem is solved if
    \begin{enumerate}
        \item[\textbf{P1}] \textbf{Convex-hull fencing:}
        $p_0(\infty)\in \co(p(\infty))$, i.e.\ the target is asymptotically
        enclosed by the convex hull spanned by the agents;
        \item[\textbf{P2}] \textbf{Collision avoidance:}
        for a given safety distance $d>0$,
        $\norm{p_{ij}(t)} > d$, $\forall\, i,j\in\NN,\ j\neq i,\ \forall\, t\ge 0$;
        \item[\textbf{P3}] \textbf{Velocity match:}
        $\lim_{t\to\infty}\norm{v_i-v_0}=0$, $\forall\, i\in\NN$.
    \end{enumerate}
    Here $v_i=\dot p_i$, $v_0=\dot p_0$, and
    \[
    \co(p) = \Bigl\{ \sum_{i=1}^N \lambda_i p_i
    : \lambda_i \ge 0,\ \sum_{i=1}^N \lambda_i = 1 \Bigr\}
    \]
    is the convex hull of the agents (area of the polygon in $\RR^2$,
    volume of the polytope in $\RR^3$).
\end{definition}

The distance from the target to the agents' convex hull is
\begin{equation}
    \dist(p,p_0)=\inf_{p^*\in\co(p)}\norm{p^*-p_0}.
    \label{eq:dist}
\end{equation}
Because $\bar p\in\co(p)$, one always has
$\dist(p,p_0)\le \norm{\bar p-p_0}=\norm{\tilde p}$; hence
$\tilde p\to 0$ implies P1. Under a baseline fencing controller
$u_i=-k_1 p_{i0}-k_2 v_{i0}+u_0+\phi_i+\psi_i$ (with the collision-avoidance
term $\phi_i$ of \S\ref{subsec:apf} and a configuration term $\psi_i$),
\eqref{eq:centroid-dyn} becomes the Hurwitz system
$\ddot{\tilde p}=-k_1\tilde p-k_2\dot{\tilde p}+\zeta$,
$\zeta=\frac1N\sum_i(\psi_i+d_i)-d_0$, which yields
$\tilde p\to 0$ only \emph{asymptotically} and with no user-specifiable
transient. Prescribing the convergence behavior is the subject of the next
subsection.

\subsection{Prescribed-performance preliminaries}
\label{subsec:ppc}

Prescribed-performance control lets the user specify, \emph{a priori}, a
single scalar performance function $\rho(t)>0$ so that the Euclidean norm of
the enclosure error satisfies
\begin{equation}
    \norm{\tilde p(t)} < \rho(t), \qquad \forall\, t\ge 0.
    \label{eq:ppc-constraint}
\end{equation}
The function $\rho(t)$ is smooth, strictly positive, decreasing,
with $\rho(0)>\norm{\tilde p(0)}$ and
$\lim_{t\to\infty}\rho(t)=\rho_\infty>0$ (taking $\rho_\infty=0$ gives
asymptotic convergence to zero). It simultaneously prescribes:
\begin{itemize}
    \item the maximum overshoot: $\norm{\tilde p(t)}<\rho(0)$;
    \item the convergence rate (via the decay rate of $\rho$);
    \item the steady-state error bound $\rho_\infty$.
\end{itemize}
Since $\dist(p,p_0)\le\norm{\tilde p(t)}$, the prescribed 2-norm envelope on
$\tilde p$ directly bounds the enclosure distance $\dist(p,p_0)$.

\textbf{Asymptotic (exponential) performance function:}
\begin{equation}
    \rho(t)=(\rho_0-\rho_\infty)e^{-\ell t}+\rho_\infty,
    \qquad \rho_0>\norm{\tilde p(0)},\ \ell>0.
    \label{eq:rho-exp}
\end{equation}

\textbf{Prescribed-time performance function:} to force the error inside
the envelope by a user-specified finite time $T$, use a finite-time
reaching profile
\begin{equation}
    \rho(t)=
    \begin{cases}
    \bar\rho + (\rho_0-\bar\rho)\bigl(1-\tfrac{t}{T}\bigr)^r, & 0\le t<T,\\[4pt]
    \bar\rho, & t\ge T,
    \end{cases}
    \qquad r\ge 2,
    \label{eq:rho-pt}
\end{equation}
which reaches the target value $\bar\rho$ exactly at $t=T$.

To convert the constrained error \eqref{eq:ppc-constraint} into an
\emph{unconstrained} variable, introduce the smooth, strictly increasing map
$T(\cdot)=\tanh(\cdot)\in(-1,1)$, let the (scalar) enclosure norm be
$e(t)=\norm{\tilde p(t)}$, and define the normalized error $z=e/\rho$ via
\[
\begin{aligned}
    &z = T(\xi)\quad\Longleftrightarrow
    \quad\\
    &\xi = T^{-1}(z)
    = \operatorname{artanh}\!\left(\frac{e}{\rho}\right)
    = \frac{1}{2}\ln\frac{1+e/\rho}{1-e/\rho}.
\end{aligned}
\]
As long as $\xi(t)$ stays finite, $|z(t)|=|T(\xi)|<1$, so
\eqref{eq:ppc-constraint} holds automatically. The norm derivative (for
$\tilde p\neq 0$) is $\dot e = \frac{\tilde p^T\dot{\tilde p}}{\norm{\tilde p}}$;
the control law derived below is continuously extensible to $\tilde p=0$
(the apparent singularity is removable, see \S\ref{subsec:ppc-design}). With
$\zeta=1/(1-z^2)>0$,
\begin{equation}
    \dot\xi
    = \frac{\zeta}{\rho}
      \bigl(\dot e - z\dot\rho\bigr),
    \qquad z=\frac{e}{\rho}.
    \label{eq:xi-dyn}
\end{equation}
Thus the fencing problem reduces to designing a controller that keeps the
single transformed error $\xi$ bounded (and drives it to zero)---the
standard PPC backstepping task for the double-integrator centroid dynamics
\eqref{eq:centroid-dyn}.

\subsection{PPC fencing problem statement}

\begin{definition}[PPC fencing property]
    \label{def:ppc-fencing}
    The fencing is achieved with \emph{prescribed performance} if, in
    addition to P1--P3, the enclosure error satisfies
    \begin{enumerate}
        \item[\textbf{P4}] \textbf{Prescribed-performance enclosure:}
        $\norm{\tilde p(t)} \le \rho(t)$ for all $t\ge 0$, where
        $\rho(t)$ is a user-specified performance function of the form
        \eqref{eq:rho-exp} or \eqref{eq:rho-pt}.
    \end{enumerate}
    Property P4 implies P1 (since $\rho(t)\to\rho_\infty$, and
    $\rho_\infty=0$ gives exact enclosure); the novelty is that the
    \emph{transient} behavior of the enclosure is now explicitly bounded by
    the user-chosen 2-norm envelope.
\end{definition}

\begin{problem}[Prescribed-performance fencing control]
    \label{prob:main}
    Given the double-integrator dynamics
    \eqref{eq:agent}--\eqref{eq:target} and a prescribed performance function
    $\rho(t)$, design a \emph{distributed} control law $u_i$ using
    only local information such that the closed loop satisfies:
    \begin{enumerate}
        \item[\textbf{P4}] prescribed-performance enclosure:
        $\norm{\tilde p(t)} < \rho(t)$, $\forall\, t\ge 0$;
        \item[\textbf{P2}] collision avoidance (hard ``fence''):
        $\norm{p_{ij}(t)}>d$ for all $i\neq j,\ t\ge 0$;
        \item[\textbf{P3}] velocity match:
        $\lim_{t\to\infty}\norm{v_i-v_0}=0$.
    \end{enumerate}
    P2 is enforced as a hard constraint via the artificial-potential barrier
    below, while P4 prescribes the soft 2-norm performance envelope on the
    enclosure error.
\end{problem}

\subsection{Artificial-potential (collision-avoidance) fence}
\label{subsec:apf}

The repulsive artificial-potential (APF) term used to enforce P2 is
\[
\phi_i^r = \sum_{j \in \NN,\ j\neq i} \alpha_r(\norm{p_{ij}})\,
\frac{p_{ij}}{\norm{p_{ij}}}
= \sum_{j \in \NN_i(t)} \alpha_r(\norm{p_{ij}})\,
\frac{p_{ij}}{\norm{p_{ij}}},
\]
where
\[
\NN_i(t)=\{j\in\NN\setminus\{i\}:\norm{p_{ij}(t)}<\mu\}
\]
is the active neighbor set, $d>0$ is the safety distance and $\mu>d$ the
activation radius. The continuous function
$\alpha_r:(d,\infty)\to[0,\infty)$ satisfies
\begin{enumerate}
    \item[A1.] $\displaystyle\lim_{s\to d^+}\int_s^\mu \alpha_r(\sigma)\,\mathrm{d}\sigma = +\infty$;
    \item[A2.] $\alpha_r(s)=0$, $\forall\, s\in[\mu,\infty)$.
\end{enumerate}
We set $\phi_i=\phi_i^r$. Since the pairwise terms are anti-symmetric,
$\sum_{i=1}^N\phi_i=0$ on the safe set
$\{p:\norm{p_{ij}}>d,\ i\ne j\}$, so the collision-avoidance potential does
not bias the swarm centroid and therefore does not interfere with the PPC
enclosure objective.
